\documentclass{article}

% ============================================================
% MA 213 Lab Activity Template
% Student-facing worksheet template for lab development.
% ============================================================

\usepackage[margin=1in]{geometry}
\usepackage{xcolor}
\usepackage{parskip}
\usepackage{array}
\usepackage{enumitem}
\usepackage{listings}

\definecolor{codegray}{gray}{0.97}
\definecolor{huddleblue}{RGB}{214,234,255}
\definecolor{predictionyellow}{RGB}{255,242,200}
\definecolor{debatepink}{RGB}{255,225,225}

\lstset{
  basicstyle=\ttfamily\small,
  columns=fullflexible,
  frame=single,
  framerule=0.4pt,
  backgroundcolor=\color{codegray},
  breaklines=true,
  showstringspaces=false,
  keepspaces=true,
  xleftmargin=0pt,
  xrightmargin=0pt
}

\newcommand{\coursenumber}{MA 213}
\newcommand{\weekplaceholder}{1}
\newcommand{\labplaceholder}{1}
\newcommand{\labtitle}{Welcome to Lab: Getting Set Up in R}

\newcommand{\nameblank}{\rule{1.75in}{0.4pt}}
\newcommand{\idblank}{\rule{1.55in}{0.4pt}}
\newcommand{\shortblank}{\rule{1.6in}{0.4pt}}
\newcommand{\longblank}{\rule{0.93\linewidth}{0.4pt}}

\newcommand{\responsebox}[2]{
  \noindent\fbox{%
    \begin{minipage}{0.95\linewidth}
      \textbf{#1} #2
      \vspace{0.55in}
    \end{minipage}
  }
}

\newcommand{\linedbox}[1]{
  \noindent\fbox{%
    \begin{minipage}{0.95\linewidth}
      #1
      \vspace{0.18in}

      \longblank

      \vspace{0.18in}
      \longblank

      \vspace{0.18in}
      \longblank
    \end{minipage}
  }
}

\newcommand{\callout}[3]{
  \noindent\colorbox{#1}{%
    \makebox[\dimexpr0.95\linewidth-2\fboxsep\relax][l]{\textbf{#2}}%
  }

  \noindent\fbox{%
    \begin{minipage}{0.93\linewidth}
      #3
    \end{minipage}
  }
}

\begin{document}

\begin{center}
  {\LARGE \coursenumber{} Week \weekplaceholder{}: Lab \labplaceholder{}}\\
  {\large \labtitle{}}
\end{center}

\begin{enumerate}[label=\arabic*.,leftmargin=*,itemsep=.7em]
  \item \textbf{Your name:} \nameblank \hfill \textbf{BU ID:} \idblank
  \item \textbf{Partner's name:} \nameblank \hfill \textbf{BU ID:} \idblank
\end{enumerate}

\textbf{Big idea:} This first lab is about getting set up and learning how lab works. There is no statistics to learn today. Your goal is to leave with R running, Tutorial 0 submitted, and this worksheet handed in.

\textbf{Do not use GenAI tools for this lab.} The point is to practice your own analysis work, coding skills, and statistical reasoning.

\textbf{Work with a partner.} One worksheet per pair, with both names above.

\textbf{Submit this worksheet at the end of the lab.}

\textbf{Files used:} \texttt{sleep.csv} (and \texttt{Lab1.R} as your starter script)

\section*{Part 1. Introduce Yourself}

Before opening R, introduce yourself to your partner. Take a minute each, then fill in the boxes below.

\responsebox{Your partner's name and major:}{}

\responsebox{Warm-up:}{What is one thing you want your partner to know about you? (Write your partner's answer here.)}

\section*{Part 2. Is R Working?}

Open RStudio. Type the following into the \textbf{Console} and press Enter.

\begin{lstlisting}
R.version.string
\end{lstlisting}

\begin{center}
\renewcommand{\arraystretch}{2.3}
\begin{tabular}{|p{0.42\linewidth}|p{0.50\linewidth}|}
\hline
\textbf{Question} & \textbf{Your answer} \\
\hline
Did RStudio open successfully? (yes / no) & \\
\hline
What version of R does it print? & \\
\hline
\end{tabular}
\end{center}

\callout{huddleblue}{HUDDLE UP}{Remember that \textbf{R} and \textbf{RStudio} are two different programs. R is the engine that does the computing; RStudio is the workspace you type into. If RStudio opens but nothing works, R is probably missing. Install R first, then RStudio. If you are stuck, raise your hand now rather than waiting.}

\newpage

\section*{Part 3. Your First Data Import}

Download \texttt{sleep.csv} and \texttt{Lab1.R} from Blackboard and put them \textbf{in the same folder}. Open \texttt{Lab1.R} in RStudio and use it as your starter script.

\begin{lstlisting}
# Check where R is currently looking for files.
getwd()

# Read in the sleep data.
sleep_data <- read.csv("sleep.csv")

# Take a first look at the top few rows.
head(sleep_data)
\end{lstlisting}

\callout{huddleblue}{HUDDLE UP}{If you get an error like ``cannot open file,'' R is looking in the wrong folder. Run \texttt{getwd()} to see where it is looking, then use \textbf{Session $\rightarrow$ Set Working Directory $\rightarrow$ To Source File Location}, or move \texttt{sleep.csv} next to your script.}

Now compute the total of the \texttt{extra} column, which records extra hours of sleep.

\begin{lstlisting}
sum(sleep_data$extra)
\end{lstlisting}

\begin{center}
\renewcommand{\arraystretch}{2.3}
\begin{tabular}{|p{0.42\linewidth}|p{0.50\linewidth}|}
\hline
\textbf{Question} & \textbf{Your answer} \\
\hline
What number does \texttt{sum(sleep\_data\$extra)} print? & \\
\hline
Which variable names does \texttt{head()} show you? & \\
\hline
\end{tabular}
\end{center}

\linedbox{\textbf{In one sentence:} what did the \texttt{\$} symbol do in \texttt{sleep\_data\$extra}?}

\section*{Part 4. Tutorial 0 and Your Hash Code}

Work through \textbf{Tutorial 0}. At the end it will generate a \textbf{hash code}.

\callout{debatepink}{IMPORTANT}{Your hash code must be submitted \textbf{to Blackboard}. Writing it on this worksheet does \textbf{not} count as submitting it. Submit it to Blackboard first, then answer the question below.}

\begin{center}
\renewcommand{\arraystretch}{2.3}
\begin{tabular}{|p{0.62\linewidth}|p{0.30\linewidth}|}
\hline
\textbf{Question} & \textbf{Your answer} \\
\hline
Did you complete Tutorial 0? (yes / no) & \\
\hline
Did you submit your hash code to Blackboard? (yes / no) & \\
\hline
Did your partner submit theirs to Blackboard? (yes / no) & \\
\hline
\end{tabular}
\end{center}

\callout{huddleblue}{BEFORE YOU LEAVE}{%
R and RStudio open without errors: \shortblank\\[.6em]
Tutorial 0 hash submitted to Blackboard: \shortblank\\[.6em]
This worksheet handed in, with both names on it: \shortblank
}

\section*{Optional Challenge}

If you and your partner finish early, import the second data set and build a simple table.

\begin{lstlisting}
burger_data <- read.csv("burger.csv")

# Count how many people chose each burger place.
table(burger_data$best_burger_place)
\end{lstlisting}

\linedbox{\textbf{What is the most visited burger place, and how can you tell from the table?}}

\end{document}
